%% ================================================================
%% SFST v30 — Neuer Supplement-Anhang: Nonperturbative Bounds
%% Titel: Appendix Q: Nonperturbative estimates and the role of
%%         QCD corrections
%%
%% Einzufügen nach Appendix P (Holographic normalization).
%% Tier-Stufen: Tier 1 für Skalenanalyse und Bounds;
%% Tier 2 für Interpretation der Kompatibilität.
%%
%% Wichtig: Das Dokument zeigt, was NICHT widerlegt wird (Hauptschlüsse)
%% UND was als offener Punkt sichtbar bleiben muss (ppb-Präzision).
%% ================================================================

\section{Appendix Q: Nonperturbative estimates and the scope of the
  $\alpha\cdot\alpha_s$ corrections}
\label{app:nonperturbative}

This appendix records conservative bounds for sea-quark,
confinement, and cross-coupling contributions of order
$\alpha\cdot\alpha_s$ to the proton--electron mass ratio.  Its goal
is not to suppress these contributions but to establish precisely
which SFST conclusions they constrain and which they leave intact.
All size estimates are stated in parts per million relative to
$m_p/m_e$ to allow direct comparison with the formula's ppm-level
corrections.

%% ----------------------------------------------------------------
\subsection{Scale inventory: what enters $m_p/m_e$}
\label{ssec:scaleInventory}

Table~\ref{tab:scales} collects the relevant contributions to
$m_p/m_e$ in decreasing order of magnitude.

\begin{table}[h]
\centering
\renewcommand{\arraystretch}{1.3}
\begin{tabular}{lrrl}
\toprule
Contribution & Size (ppm) & MeV equiv. & Source \\
\midrule
QCD confinement anomaly $H_a/m_p$ &
  $\sim 2\times10^5$ & $\sim 200$ &
  EMT trace anomaly; Yang et al.\ 2018 \\
Gluon energy $H_g/m_p$ &
  $\sim 3.4\times10^5$ & $\sim 338$ &
  Lattice QCD; Yang et al.\ 2018 \\
Strange-quark $\sigma$-term $\sigma_s/m_p$ &
  $\sim 4.4\times10^4$ & $\sim 41$ &
  Bali et al.\ 2016 \\
EM self-energy $\delta m_p^{\mathrm{EM}}/m_p$ &
  $\sim 670$ & $\sim 0.63$ &
  Cottingham sum rule \\
$\alpha\cdot\alpha_s$ cross-correction &
  $\sim 100$--$500$ & $\sim 0.1$--$0.5$ &
  ChPT estimate (§\ref{ssec:crossBound}) \\
\midrule
SFST geometric baseline shortfall &
  $18.82$ & $0.0346$ & $6\pi^5$ vs.\ CODATA \\
SFST collective correction $\alpha^2/\sqrt{8}$ &
  $18.83$ & $0.0346$ & \cref{eq:master} \\
\bottomrule
\end{tabular}
\caption{Size inventory of contributions to $m_p/m_e$, ordered by
  magnitude.  The SFST correction is the second-to-last entry.
  All ppm values are relative to the experimental ratio
  $\mu_{\exp}=1836.152673$.}
\label{tab:scales}
\end{table}

The first observation from Table~\ref{tab:scales} is structural:
\emph{the QCD contributions to $m_p$ dwarf the SFST correction by
factors of $10^3$--$10^4$}.  The SFST framework is therefore not
operating in a perturbative regime around a QCD-free baseline.

%% ----------------------------------------------------------------
\subsection{Conservative bound on the $\alpha\cdot\alpha_s$ term}
\label{ssec:crossBound}

The leading electroweak--strong cross-coupling arises from
electromagnetic dressing of the QCD self-energy diagram.  Using the
operator product expansion with $\overline{\mathrm{MS}}$ matching:
\begin{equation}
  \label{eq:crossTerm}
  \left.\frac{\delta(m_p/m_e)}{m_p/m_e}\right|_{\alpha\cdot\alpha_s}
  \approx \frac{\alpha\,\alpha_s\,C_F}{4\pi^2}
          \ln\!\frac{m_p}{\Lambda_{\mathrm{QCD}}}
  \approx \frac{(7.30\times10^{-3})(0.4)(4/3)}{4\pi^2}
          \ln\!\frac{938}{332}
  \approx 1.0\times10^{-4}.
\end{equation}
In ppm this is approximately $100$~ppm.  A more conservative upper
bound, accounting for non-logarithmic terms and for the uncertainty
in $\alpha_s$ at the proton scale, places the cross-correction in the
range
\begin{equation}
  \label{eq:crossBound}
  \left|\frac{\delta(m_p/m_e)}{m_p/m_e}\right|_{\alpha\cdot\alpha_s}
  \lesssim 500\;\mathrm{ppm}.
\end{equation}

Comparing with the SFST correction $\alpha^2/\sqrt{8}\approx 18.8$~ppm,
the cross-term is estimated to be
\begin{equation}
  \frac{|\delta_{\alpha\cdot\alpha_s}|}{\alpha^2/\sqrt{8}}
  \;\approx\; 5\text{--}27.
\end{equation}

This ratio constitutes the main obstruction to interpreting the SFST
correction as the \emph{unique} leading electromagnetic contribution
in a conventional perturbative expansion: in any such expansion,
$\alpha\cdot\alpha_s$ terms appear at an earlier order and are
numerically larger than $\alpha^2$.

%% ----------------------------------------------------------------
\subsection{What the bounds do not overturn}
\label{ssec:compatible}

Despite the size of the corrections, the bounds do not invalidate the
following claims, which remain the defensible conclusions of the paper.

\paragraph{Claim 1 (Geometric baseline, Tier~1).}
$6\pi^5 = 1836.118\ldots$ lies within $18.8$~ppm of the experimental
value.  This 18.8-ppm discrepancy is \emph{smaller than, or
comparable to}, the typical size of EM--QCD cross-corrections
($\sim 100$--$500$~ppm), which means the geometric term $6\pi^5$ is
numerically plausible as a zeroth-order approximation even within a
full QCD+EM framework.  The bounds do not rule out a scenario in which
$6\pi^5$ organises the leading geometric scale and
nonperturbative corrections accumulate at $\lesssim 100$~ppm.

\paragraph{Claim 2 (Residual structure, observation-level).}
The residual $\mu_{\exp}-6\pi^5 \approx 18.8\text{ ppm}$ is
numerically close to $\alpha^2/\sqrt{8}\approx 18.8\text{ ppm}$.
This numerical coincidence is retained as an \emph{observation}
(labelled as such in the main text).  The bounds do not
make this coincidence less striking; they do constrain the physical
interpretation (see §\ref{ssec:precisionClaim}).

\paragraph{Claim 3 (Tier~1 stability result).}
The first-order zeta-pole residue rigidity under volume-preserving
anisotropic deformations is a pure spectral-geometric theorem.  It
does not involve QCD and is therefore unaffected by the bounds.

\paragraph{Claim 4 (Tadpole cancellation, Tier~2).}
The universal vanishing of the electromagnetic tadpole
$\partial_\theta Z|_{\theta=0}=0$ (Appendix~R) follows from the
lattice antisymmetry $n_j\mapsto -n_j$ and is independent of
$\alpha_s$.  The bounds do not affect this structural result.

%% ----------------------------------------------------------------
\subsection{What the bounds constrain: the precision claim}
\label{ssec:precisionClaim}

The bounds do constrain a stronger claim that is \emph{not} made in
the current manuscript but appears in informal discussions of SFST:

\begin{quote}
\emph{``The formula $6\pi^5(1+\alpha^2/\sqrt{8})$ reproduces
$m_p/m_e$ to ppb accuracy as a genuine first-principles derivation.''}
\end{quote}

This claim is inconsistent with the size of $\alpha\cdot\alpha_s$
corrections for the following reason.  The formula matches experiment
to approximately $2$~ppb ($\approx 0.002$~ppm).  The
$\alpha\cdot\alpha_s$ correction is estimated at $\sim100$~ppm.
For the ppb-level match to be non-accidental, one of the following
must hold:
\begin{enumerate}[label=(\roman*)]
  \item The SFST framework already absorbs all QCD contributions
        into the geometric factor $6\pi^5$, and the surviving
        correction is precisely $\alpha^2/\sqrt{8}$.  This would
        require an explicit QCD cancellation mechanism, currently
        absent from the framework.
  \item The $\alpha\cdot\alpha_s$ corrections vanish by a structural
        cancellation between the proton and electron sectors.  The
        electron has no QCD interactions, so such a cancellation
        would require a specific matching between the hadronic and
        leptonic spectral observables that is not yet specified.
  \item The ppb-level agreement is accidental (consistent with the
        framework's own characterisation of the formula as a
        ``numerically suggestive observation'' rather than a derivation).
\end{enumerate}

The present manuscript adopts interpretation~(iii), which is the
only one currently supported by the available derivations.
Interpretations~(i) and~(ii) are possible routes for future work but
require additional structural inputs.

%% ----------------------------------------------------------------
\subsection{Summary table: which conclusions survive}
\label{ssec:boundsSummary}

\begin{center}
\renewcommand{\arraystretch}{1.3}
\begin{tabular}{p{0.46\textwidth}cc}
\toprule
Conclusion & Tier & Survives bounds? \\
\midrule
$6\pi^5$ lies within 20\,ppm of $\mu_{\exp}$ &
  1 (arithmetic) & \checkmark~yes \\
$\alpha^2/\sqrt{8}$ closes the residual to ppb & 
  obs.\ (not claimed as derivation) & \checkmark~as observation \\
First-order zeta-pole rigidity under anisotropy &
  1 (theorem) & \checkmark~yes \\
Universal tadpole cancellation $\partial_\theta Z|_0=0$ &
  1 (theorem) & \checkmark~yes \\
$\alpha^2/\sqrt{8}$ is the \emph{leading} EM correction &
  open & $\times$~constrained \\
ppb precision is a first-principles prediction &
  not claimed & $\times$~not supported \\
\bottomrule
\end{tabular}
\end{center}

The nonperturbative bounds therefore serve a clarifying rather than a
destructive function: they identify precisely which part of the SFST
program requires an additional QCD cancellation theorem and which
part is already protected by operator-level structural arguments.
